% Section 6.1, from lectures/L06/appendix/a_sequential_logic.md.

\section{Sekvensnät och vippor}
\label{c6:app:a}

\subsection{Från kombinatoriskt till sekventiellt}
\label{c6:a1}
Alla nät hittills har varit \textbf{kombinatoriska}: utgången beror bara på ingångarnas värden just
nu. Samma ingångar ger alltid samma utgång, och nätet har ingen aning om vad som hände för en
sekund sedan. Bromsassistenten, majoritetskretsen och segmentavkodaren är alla sådana.

Många verkliga uppgifter kräver mer. Ett transportband som startas med en knapp och stoppas med en
annan måste \textbf{minnas} att startknappen har tryckts, även efter att den släppts. En räknare
måste minnas hur långt den räknat. Ett trafikljus måste minnas vilken färg som lyser, för att veta
vilken som kommer härnäst. En mikrodator måste minnas sitt program, sina variabler och vilken
instruktion den är på.

Ett nät som minns kallas ett \textbf{sekvensnät}: utgången beror på ingångarna \textbf{och} på vad
som hänt tidigare, alltså på ett lagrat \textbf{tillstånd}. Det här avsnittet visar hur minne byggs
av vanliga grindar, och vad som byggs av minnet: register, räknare och skiftregister. Det är de
byggstenar som mikrodatorn i \secref{c6:app:b} består av.

Sekvensnät ingår inte i kursplanens lärandemål för sin egen skull. De finns med för att
mikrodatorn annars blir en svart låda: dess register, programräknare och minnen är vippor.

\subsection{Minne genom återkoppling}
\label{c6:a2}
Du har redan byggt ett sekvensnät, i kapitel~\ref{c2:ch}: reläets \textbf{självhållning}
(\secref{c2:b7}).

\bookfigure[tbp]{l06_self_hold}{Självhållning med startknappen S1, stoppknappen S2 och reläet
K1.}{c6:fig:self_hold}

Tryck på start, \code{S1}, så får reläspolen \code{K1} ström och drar. Reläets egen kontakt
\code{K1} sitter parallellt med \code{S1}, så när \code{S1} släpps går strömmen i stället genom
\code{K1}:s kontakt, och reläet håller sig självt draget. Först när stopp, \code{S2}, trycks bryts
strömmen, reläet släpper, och kretsen är tillbaka i utgångsläget.

Samma knapptryckning, \code{S1} släppt och \code{S2} släppt, ger alltså olika resultat beroende på
vad som hänt innan: lampan lyser om start senast trycktes, och är släckt om stopp senast trycktes.
Kretsen minns.

Nyckeln är \textbf{återkopplingen}: reläets utgång, kontakten, är kopplad tillbaka till dess egen
ingång. Samma idé fungerar med grindar.

\subsection{SR-låset}
\label{c6:a3}
Koppla två NOR-grindar så att varje grinds utgång matar den andra grindens ingång:

\bookfigure{l06_sr_latch}{SR-lås av två korskopplade NOR-grindar.}{c6:fig:sr_latch}

Det här är ett \textbf{SR-lås} (\emph{SR latch}), den enklaste minnescellen som finns. Den har två
ingångar:
\begin{itemize}
  \item \textbf{S} (\emph{set}, ettställ): gör \code{Q = 1};
  \item \textbf{R} (\emph{reset}, nollställ): gör \code{Q = 0}.
\end{itemize}
Utgången \code{Q'} är alltid inversen av \code{Q}, utom i ett fall (nedan).

\begin{narrowtable}{@{}ccll@{}}
  \thead{S} & \thead{R} & \thead{Q efteråt} & \thead{Betydelse} \\
  \midrule
  0 & 0 & samma som före & \textbf{minns} \\
  1 & 0 & 1 & ettställ \\
  0 & 1 & 0 & nollställ \\
  1 & 1 & ogiltigt & båda utgångarna 0 \\
\end{narrowtable}

Följ signalerna en gång för att se varför det minns. Anta att \code{Q = 1}, så \code{Q' = 0}, och
att \code{S = R = 0}. Den övre grinden får \code{R = 0} och \code{Q' = 0}, och en NOR av två nollor
är 1: \code{Q} förblir 1. Den nedre grinden får \code{Q = 1} och \code{S = 0}, och en NOR med en
etta in är 0: \code{Q'} förblir 0. Tillståndet håller sig självt, precis som reläet. Ett kort tryck
på \code{R} gör \code{Q = 0}, och sedan håller sig det tillståndet lika bra.

Jämför med självhållningen: \code{S1} är \textbf{S}, \code{S2} är \textbf{R}, och reläet är låset.
Skillnaden är fallet \code{S = R = 1}. I reläkretsen vinner stopp, eftersom \code{S2} sitter i
serie med allt; i NOR-låset blir båda utgångarna 0, och vad som händer när båda släpps samtidigt
går inte att förutsäga. Därför är kombinationen förbjuden, och därför används SR-låset sällan
direkt.

Bygg låset i CircuitVerse (\url{https://circuitverse.org/simulator}) och prova: tryck \code{S},
släpp, tryck \code{R}, släpp. Lägg märke till att utgången inte går att förutsäga från ingångarna
ensamma. Du måste veta vad som hände innan.

\subsection{D-låset}
\label{c6:a4}
\textbf{D-låset} tar bort den förbjudna kombinationen. I stället för \code{S} och \code{R} har det
en dataingång \code{D}, värdet som ska lagras, och en ingång \code{en} (\emph{enable}) som
bestämmer när låset lyssnar:

\bookfigure{l06_d_latch}{D-lås: en inverterare och två AND-grindar framför ett
SR-lås.}{c6:fig:d_latch}

De två AND-grindarna gör \code{S = D · en} och \code{R = D' · en}. När \code{en = 0} är både
\code{S} och \code{R} noll, och låset minns. När \code{en = 1} blir antingen \code{S} eller
\code{R} en etta, beroende på \code{D}, men aldrig båda.

\begin{narrowtable}{@{}cl@{}}
  \thead{en} & \thead{Beteende} \\
  \midrule
  1 & \textbf{genomsläppligt}: \code{Q} följer \code{D} direkt \\
  0 & \textbf{låst}: \code{Q} behåller det värde det hade när \code{en} gick till 0 \\
\end{narrowtable}

D-låset är \textbf{nivåkänsligt}: det släpper igenom allt som händer på \code{D} så länge
\code{en} är 1. Det är också dess svaghet. En störning på \code{D} medan \code{en = 1} går rakt
igenom till \code{Q}, och i ett system med tusentals minnesceller är det svårt att garantera att
varje \code{D} är stabilt under hela den tid \code{en} är 1.

\subsection{D-vippan och klockan}
\label{c6:a5}
\textbf{D-vippan} (\emph{D flip-flop}) löser problemet. Den lagrar också en bit, men den läser
\code{D} bara i ett enda \textbf{ögonblick}: när en \textbf{klocksignal} går från 0 till 1, den
\textbf{stigande flanken}. Resten av tiden, oavsett vad \code{D} gör, behåller \code{Q} sitt
värde.

Invändigt är en D-vippa två D-lås i rad, där det ena är öppet när klockan är låg och det andra när
den är hög. Nettoeffekten är att bara värdet i själva ögonblicket när klockan stiger kommer
igenom.

Kör ett lås och en vippa med samma klocka och samma \code{D}, så syns skillnaden:

\bookfigure[tbp]{l06_latch_vs_flipflop}{Ett D-lås och en D-vippa med samma klocka och samma
\code{D}.}{c6:fig:latch_vs_flipflop}

Låset ändrar sig flera gånger, varje gång \code{D} ändras medan klockan är hög. Vippan ändrar sig
bara vid de stigande flankerna, och bara om \code{D} då har ett annat värde än det \code{Q} redan
har. En störning mellan två flanker syns inte alls:

\bookfigure{l06_dff_timing}{En kort störning på \code{D} mellan två stigande
flanker.}{c6:fig:dff_timing}

\textbf{Klockan} är en fyrkantvåg som alla vippor i ett system delar. Två tal beskriver den:
\textbf{perioden} \code{T}, tiden för en hel svängning, och \textbf{frekvensen} \code{f = 1/T}. En
Arduino Uno har en klocka på \textbf{16~MHz}, sexton miljoner svängningar i sekunden, så perioden
är
\[
  T = \frac{1}{16\,000\,000\ \text{Hz}} = 62{,}5\ \text{ns}
\]

Varje vippa i mikrodatorn uppdateras vid samma flank. Det gör att allt i datorn tar ett steg i
taget, i takt, och det är därför det går att säga att en instruktion tar ett bestämt, helt antal
\textbf{klockcykler}. Det talet kommer tillbaka i kapitel~\ref{c9:ch}, där tidsfördröjningar
räknas ut i klockcykler.

I scheman ritas en D-vippa som en ruta med ingången \code{D}, utgången \code{Q} och en liten
triangel vid klockingången. Triangeln betyder \enquote{flanktriggad}.

\subsection{Register}
\label{c6:a6}
En vippa lagrar en bit. Ett \textbf{register} lagrar flera: det är N vippor sida vid sida, som delar
samma klocka.

\bookfigure[tbp]{l06_register}{Ett 4-bitars register av fyra D-vippor med gemensam
klocka.}{c6:fig:register}

Vid varje stigande flank kopieras alla ingångarna \code{D3}--\code{D0} till utgångarna
\code{Q3}--\code{Q0} samtidigt, och resten av tiden håller registret sitt värde.

Mikrodatorn i resten av boken är full av register. De 32 \textbf{arbetsregistren}
\code{r0}--\code{r31}, där programmet räknar, är 8-bitars register: åtta vippor vardera.
\textbf{Statusregistret} \code{SREG}, där räkneenheten sparar sina flaggor, är ett register. Och
\textbf{utporten} \code{PORTB}, som tänder lysdioderna i Labb~3, är ett register vars utgångar är
kopplade direkt till kretsens stift: skriv en etta i en bit, så blir stiftet 5~V.

\subsection{Räknare}
\label{c6:a7}
Koppla en vippa så att den \textbf{växlar} vid varje flank: \code{D} får inversen av \code{Q}. Då
blir \code{Q} en fyrkantvåg med halva klockans frekvens. Kedja flera sådana, så att varje vippa
växlar när den föregående går från 1 till 0, så får du en \textbf{räknare}:

\bookfigure{l06_counter_timing}{Tidsdiagram för en 3-bitars räknare.}{c6:fig:counter_timing}

Läs \code{Q2 Q1 Q0} som ett binärt tal efter varje flank: 001, 010, 011, 100 och så vidare upp till
111, och sedan 000 igen. Tre vippor räknar till 7; \code{n} vippor räknar från 0 till
\code{2^n - 1} och börjar sedan om. Att räknaren slår om från 111 till 000 är ingen särskild logik;
det är vad som händer när den fjärde biten inte finns. Samma sak kallas \textbf{rundgång} i
kapitel~\ref{c8:ch}, när ett 8-bitars register räknar förbi 255.

Två saker att lägga märke till:
\begin{itemize}
  \item \textbf{Varje bit har halva frekvensen av biten före.} \code{Q0} växlar vid varje flank,
    \code{Q1} vid varannan, \code{Q2} vid var fjärde. En räknare är alltså också en
    frekvensdelare.
  \item \textbf{Mikrodatorns programräknare är en räknare.} Den håller adressen till nästa
    instruktion och räknar upp ett steg för varje instruktion som hämtas (\secref{c6:app:b}).
\end{itemize}

\subsection{Skiftregister}
\label{c6:a8}
Koppla vipporna i en kedja, så att varje vippas \code{Q} matar nästa vippas \code{D}:

\bookfigure{l06_shift_register}{Ett 4-bitars skiftregister: fyra D-vippor i
kedja.}{c6:fig:shift_register}

Vid varje flank tar varje vippa över värdet från vippan före, och den första tar in en ny bit.
Innehållet \textbf{flyttas}, skiftas, ett steg per klockflank. Det kallas ett
\textbf{skiftregister}.

\begin{narrowtable}{@{}cccccc@{}}
  \thead{Efter flank} & \thead{\code{in}} & \thead{Q0} & \thead{Q1} & \thead{Q2} & \thead{Q3} \\
  \midrule
  start & & 0 & 0 & 0 & 0 \\
  1 & 1 & 1 & 0 & 0 & 0 \\
  2 & 0 & 0 & 1 & 0 & 0 \\
  3 & 0 & 0 & 0 & 1 & 0 \\
  4 & 0 & 0 & 0 & 0 & 1 \\
\end{narrowtable}

En enda etta vandrar genom registret, som ett rinnande ljus. Skiftregister används för att skicka
data en bit i taget över en ledning (seriell överföring), och de har en direkt motsvarighet i
mikrodatorn: instruktionen \code{lsl} flyttar alla bitar i ett register ett steg åt vänster, vilket
är samma sak som att multiplicera med 2. Den kommer i kapitel~\ref{c8:ch}. Att figuren flyttar åt
höger är bara ritsättet: skrivs innehållet som ett binärt tal, \code{Q3 Q2 Q1 Q0}, flyttas varje bit
till en mer signifikant plats, alltså åt vänster i talet, precis som med \code{lsl}.
